\documentclass[10pt,a4paper]{amsart}
\usepackage{amsmath,amssymb,amsthm,graphicx}
\usepackage{verbatim}
\newenvironment{reference}{\comment}{\endcomment}
%\newenvironment{reference}{}{}
\newenvironment{slogan}{\comment}{\endcomment}
\newenvironment{history}{\comment}{\endcomment}

% For commutative diagrams we use Xy-pic
\usepackage[all]{xy}

% We use 2cell for 2-commutative diagrams.
\xyoption{2cell}
% Theorem environments.
%
\theoremstyle{plain}
\newtheorem{theorem}[subsection]{Theorem}
\newtheorem{proposition}[subsection]{Proposition}
\newtheorem{lemma}[subsection]{Lemma}

\theoremstyle{definition}
\newtheorem{definition}[subsection]{Definition}
\newtheorem{example}[subsection]{Example}
\newtheorem{exercise}[subsection]{Exercise}
\newtheorem{situation}[subsection]{Situation}

\theoremstyle{remark}
\newtheorem{remark}[subsection]{Remark}
\newtheorem{remarks}[subsection]{Remarks}

\numberwithin{equation}{subsection}

% Macros
%
\def\lim{\mathop{\mathrm{lim}}\nolimits}
\def\colim{\mathop{\mathrm{colim}}\nolimits}
\def\Spec{\mathop{\mathrm{Spec}}}
\def\Hom{\mathop{\mathrm{Hom}}\nolimits}
\def\Ext{\mathop{\mathrm{Ext}}\nolimits}
\def\SheafHom{\mathop{\mathcal{H}\!\mathit{om}}\nolimits}
\def\SheafExt{\mathop{\mathcal{E}\!\mathit{xt}}\nolimits}
\def\Sch{\mathit{Sch}}
\def\Mor{\mathop{\mathrm{Mor}}\nolimits}
\def\Ob{\mathop{\mathrm{Ob}}\nolimits}
\def\Sh{\mathop{\mathit{Sh}}\nolimits}
\def\NL{\mathop{N\!L}\nolimits}
\def\CH{\mathop{\mathrm{CH}}\nolimits}
\def\proetale{{pro\text{-}\acute{e}tale}}
\def\etale{{\acute{e}tale}}
\def\QCoh{\mathit{QCoh}}
\def\Ker{\mathop{\mathrm{Ker}}}
\def\Im{\mathop{\mathrm{Im}}}
\def\Coker{\mathop{\mathrm{Coker}}}
\def\Coim{\mathop{\mathrm{Coim}}}

% Boxtimes
%
\DeclareMathSymbol{\boxtimes}{\mathbin}{AMSa}{"02}

%
% Macros for moduli stacks/spaces
%
\def\QCohstack{\mathcal{QC}\!\mathit{oh}}
\def\Cohstack{\mathcal{C}\!\mathit{oh}}
\def\Spacesstack{\mathcal{S}\!\mathit{paces}}
\def\Quotfunctor{\mathrm{Quot}}
\def\Hilbfunctor{\mathrm{Hilb}}
\def\Curvesstack{\mathcal{C}\!\mathit{urves}}
\def\Polarizedstack{\mathcal{P}\!\mathit{olarized}}
\def\Complexesstack{\mathcal{C}\!\mathit{omplexes}}
% \Pic is the operator that assigns to X its picard group, usage \Pic(X)
% \Picardstack_{X/B} denotes the Picard stack of X over B
% \Picardfunctor_{X/B} denotes the Picard functor of X over B
\def\Pic{\mathop{\mathrm{Pic}}\nolimits}
\def\Picardstack{\mathcal{P}\!\mathit{ic}}
\def\Picardfunctor{\mathrm{Pic}}
\def\Deformationcategory{\mathcal{D}\!\mathit{ef}}
\usepackage{cleveref}
\crefname{originallemma}{original lemma}{original lemmas}
\Crefname{originallemma}{Original lemma}{Original lemmas}
\makeatletter
\newlabel{categories-lemma-morphisms-left-localization}{{04VE}{source}{}{sourcepointer-1}{}}
\newlabel{categories-lemma-morphisms-left-localization@cref}{{[originallemma][0][]04VE}{[1][0][]source}{}{}{}}
\makeatother
\crefname{originalequation}{original equation}{original equations}
\Crefname{originalequation}{Original equation}{Original equations}
\makeatletter
\newlabel{categories-equation-left-localization-morphisms-colimit}{{05Q1}{source}{}{sourcepointer-2}{}}
\newlabel{categories-equation-left-localization-morphisms-colimit@cref}{{[originalequation][0][]05Q1}{[1][0][]source}{}{}{}}
\makeatother
\crefname{originalremark}{original remark}{original remarks}
\Crefname{originalremark}{Original remark}{Original remarks}
\makeatletter
\newlabel{categories-remark-left-localization-morphisms-colimit}{{05Q0}{source}{}{sourcepointer-3}{}}
\newlabel{categories-remark-left-localization-morphisms-colimit@cref}{{[originalremark][0][]05Q0}{[1][0][]source}{}{}{}}
\makeatother
\title{The unique preadditive structure making the localization functor additive}
\author{The Stacks Project Authors}
\date{stacks-project:homology.tex:local-snapshot-20261012; GNU FDL 1.2 or later}
\begin{document}
\maketitle
\noindent\textit{Editorial note.} All delivered mathematical text is the author's own, copied from the pinned Stacks Project chapter file homology.tex; no statement or proof step was rewritten, reordered or removed.\par\medskip
\noindent\textit{Editorial note.} The proof is the author's complete argument for one of the two cases: the author states 'We will prove this in the case S is a left multiplicative system. The case where S is a right multiplicative system is dual.', and the duality sentence closes a case that is the exact dual of the case written out in full inside the counted span, so it is the author's own routine duality remark and is preserved verbatim. Every step of the written-out case is inside the counted span and no closing sentence is a placeholder.\par\medskip
\noindent\textit{Editorial note.} Every labelled reference occurring inside the counted spans is carried in assembly.external\_references as a non-load-bearing navigation pointer (same-chapter ones at their definition site, cross-chapter ones at the invocation inside the counted span), so no load-bearing mathematics is replaced by a stub. No bibliography citation occurs inside the counted spans.\par\medskip
\noindent\textit{Editorial note.} The same-name check for this card scanned every staged Stacks chapter of this repository -- all sources/stacks-*/pinned-original-source.tex files, the card's own chapter excluded -- and the labels found are recorded in dedup.same\_name\_cross\_chapter on the card.\par\medskip
\noindent\textit{Editorial note.} The preamble is the chapter's own native preamble excerpt already used for this source: the comment-type reference, slogan and history environments (native preamble lines 8-13), the Xy-pic setup (native preamble lines 14-18) and the native theorem declarations with the mathematical macros (native preamble lines 154-217, the actual last line of the file).\par\medskip
\noindent\textit{Allowed standard prerequisites (editorial).}\begin{itemize}
\item Fraction representations of morphisms of a left localization. Categories, Lemma categories-lemma-morphisms-left-localization. Lets every morphism of the localization be written as s\textasciicircum{}\{-1\} f with a common denominator, which is how two morphisms are represented at the beginning of the proof.
\item The filtered-colimit description of the morphism sets of a left localization. Categories, Equation categories-equation-left-localization-morphisms-colimit and Remark categories-remark-left-localization-morphisms-colimit. Supplies the enlargement step and the filteredness of the category Y'/S used to prove that the proposed addition is well defined and that the compositions agree.
\item The left multiplicative system axioms. Categories, the definition of a left multiplicative system (axioms LMS1-LMS3) used throughout the localization section. Supply the morphisms s', s\_1, s\_2 and a, b chosen in the two diagrams of the proof.
\end{itemize}\medskip
\ifdefined\hypertarget\hypertarget{sourcepointer-1}{}\fi
\noindent\textit{Source reference.} Original lemma 04VE: sources/stacks-categories/pinned-original-source.tex, source lines 4485--4485. Navigation pointer invoked at proof line 1984; the counted span contains the pointer only and the cited result is not reproduced inside it, so the pointer is registered as a non-load-bearing navigation stub in the established form for Stacks material. No bibliography citation is involved.\par\medskip
\ifdefined\hypertarget\hypertarget{sourcepointer-2}{}\fi
\noindent\textit{Source reference.} Original equation 05Q1: sources/stacks-categories/pinned-original-source.tex, source lines 4614--4614. Navigation pointer invoked at proof line 2006; the counted span contains the pointer only and the cited result is not reproduced inside it, so the pointer is registered as a non-load-bearing navigation stub in the established form for Stacks material. No bibliography citation is involved.\par\medskip
\ifdefined\hypertarget\hypertarget{sourcepointer-3}{}\fi
\noindent\textit{Source reference.} Original remark 05Q0: sources/stacks-categories/pinned-original-source.tex, source lines 4580--4580. Navigation pointer invoked at proof line 2011; the counted span contains the pointer only and the cited result is not reproduced inside it, so the pointer is registered as a non-load-bearing navigation stub in the established form for Stacks material. No bibliography citation is involved.\par\medskip
\begin{lemma}
\label{lemma-localization-preadditive}
Let $\mathcal{C}$ be a preadditive category.
Let $S$ be a left or right multiplicative system.
There exists a unique preadditive structure on
$S^{-1}\mathcal{C}$ such that the localization functor
$Q : \mathcal{C} \to S^{-1}\mathcal{C}$ is additive.
\end{lemma}

\begin{proof}
We will prove this in the case $S$ is a left multiplicative system.
The case where $S$ is a right multiplicative system is dual.
Suppose that $X, Y$ are objects of $\mathcal{C}$ and that
$\alpha, \beta : X \to Y$ are morphisms in $S^{-1}\mathcal{C}$. According to
Categories, Lemma \ref{categories-lemma-morphisms-left-localization}
we may represent these by pairs $s^{-1}f, s^{-1}g$ with common denominator $s$.
Suppose that we have a pre-additive structure on $S^{-1}\mathcal{C}$ and
that $Q$ is additive. Then
\begin{align*}
s^{-1}f+s^{-1}g
&=Q(s)^{-1}\circ Q(f)+Q(s)^{-1}\circ Q(g)\\
&=Q(s)^{-1}\circ(Q(f)+Q(g))\\
&=Q(s)^{-1}\circ Q(f+g)\\
&=s^{-1}(f+g).
\end{align*}
Hence, the pre-additive structure on $S^{-1}\mathcal{C}$ is unique if it
exists. Correspondingly, we define $\alpha + \beta$ to be the
equivalence class of $s^{-1}(f + g)$ and in the rest of the proof we show
that this is well defined and that composition is bilinear.
Once this is done it is clear that $Q$ is an additive functor.

\medskip\noindent
Let us show construction above is well defined.
An abstract way of saying this is that filtered colimits of
abelian groups agree with filtered colimits of sets and to use
Categories,
Equation (\ref{categories-equation-left-localization-morphisms-colimit}).
We can work this out in a bit more detail as follows.
Say $s : Y \to Y_1$ and $f, g : X \to Y_1$. Suppose we have a second
representation of $\alpha, \beta$ as $(s')^{-1}f', (s')^{-1}g'$ with
$s' : Y \to Y_2$ and $f', g' : X \to Y_2$. By
Categories, Remark \ref{categories-remark-left-localization-morphisms-colimit}
we can find a morphism $s_3 : Y \to Y_3$ and morphisms
$a_1 : Y_1 \to Y_3$, $a_2 : Y_2 \to Y_3$ such that
$a_1 \circ s = s_3 = a_2 \circ s'$ and also
$a_1 \circ f = a_2 \circ f'$ and $a_1 \circ g = a_2 \circ g'$.
Hence we see that $s^{-1}(f + g)$ is equivalent to
\begin{align*}
s_3^{-1}(a_1 \circ (f + g)) & =
s_3^{-1}(a_1 \circ f + a_1 \circ g) \\
& = s_3^{-1}(a_2 \circ f' + a_2 \circ g') \\
& = s_3^{-1}(a_2 \circ (f' + g'))
\end{align*}
which is equivalent to $(s')^{-1}(f' + g')$.

\medskip\noindent
Fix $s : Y \to Y'$ and $f, g : X \to Y'$ with
$\alpha = s^{-1}f$ and $\beta = s^{-1}g$ as morphisms $X \to Y$
in $S^{-1}\mathcal{C}$.
To show that composition is bilinear first consider the case of a
morphism $\gamma : Y \to Z$ in $S^{-1}\mathcal{C}$. Say $\gamma = t^{-1}h$
for some $h : Y \to Z'$ and $t : Z \to Z'$ in $S$. Using LMS2 we
choose morphisms $a : Y' \to Z''$ and $t' : Z' \to Z''$ in $S$ such
that $a \circ s = t' \circ h$. Picture
$$
\xymatrix{
& & Z \ar[d]^t \\
& Y \ar[r]^h \ar[d]^s & Z' \ar[d]^{t'} \\
X \ar[r]^{f, g} & Y' \ar[r]^a & Z''
}
$$
Then
$\gamma \circ \alpha = (t' \circ t)^{-1}(a \circ f)$ and
$\gamma \circ \beta = (t' \circ t)^{-1}(a \circ g)$.
Hence we see that $\gamma \circ (\alpha + \beta)$ is represented
by $(t' \circ t)^{-1}(a \circ (f + g)) =
(t' \circ t)^{-1}(a \circ f + a \circ g)$ which represents
$\gamma \circ \alpha + \gamma \circ \beta$.

\medskip\noindent
Finally, assume that $\delta : W \to X$ is another morphism of
$S^{-1}\mathcal{C}$. Say $\delta = r^{-1}i$ for some
$i : W \to X'$ and $r : X \to X'$ in $S$. We claim that we can find
a morphism $s' : Y' \to Y''$ in $S$ and morphisms $a'', b'' : X' \to Y''$
such that the following diagram commutes
$$
\xymatrix{
& & & Y \ar[d]^s \\
& X \ar[rr]^{f, g, f + g} \ar[d]^r & & Y' \ar[d]^{s'} \\
W \ar[r]^i & X' \ar[rr]^{a'', b'', a'' + b''} & & Y''
}
$$
Namely, using LMS2 we can first choose
$s_1 : Y' \to Y_1$, $s_2 : Y' \to Y_2$ in $S$ and
$a : X' \to Y_1$, $b : X' \to Y_2$ such that
$a \circ r = s_1 \circ f$ and $b \circ r = s_2 \circ f$.
Then using that the category $Y'/S$ is filtered (see
Categories, Remark \ref{categories-remark-left-localization-morphisms-colimit}),
we can
find a $s' : Y' \to Y''$ and morphisms $a' : Y_1 \to Y''$, $b' : Y_2 \to Y''$
such that $s' = a' \circ s_1$ and $s' = b' \circ s_2$. Setting
$a'' = a' \circ a$ and $b'' = b' \circ b$ works.
At this point we see that the compositions
$\alpha \circ \delta$ and $\beta \circ \delta$ are represented by
$(s' \circ s)^{-1}(a'' \circ i)$ and $(s' \circ s)^{-1}(b'' \circ i)$.
Hence $\alpha \circ \delta + \beta \circ \delta$ is represented
by $(s' \circ s)^{-1}(a'' \circ i + b'' \circ i) =
(s' \circ s)^{-1}((a'' + b'') \circ i)$
which by the diagram again is a representative
of $(\alpha + \beta) \circ \delta$.
\end{proof}

\end{document}
